\documentclass[12pt,a4paper]{article}
\usepackage[utf8]{inputenc}
\usepackage[spanish]{babel}
\usepackage{graphicx}
\usepackage{listings}
\usepackage{xcolor}
\usepackage{hyperref}
\usepackage{amsmath}
\usepackage{geometry}
\geometry{margin=2.5cm}

% Configuración de listings para código MIPS
\lstdefinestyle{mipsstyle}{
    language={[MIPS]Assembler},
    basicstyle=\ttfamily\small,
    keywordstyle=\color{blue}\bfseries,
    commentstyle=\color{green!60!black},
    stringstyle=\color{red},
    numbers=left,
    numberstyle=\tiny,
    numbersep=5pt,
    frame=single,
    breaklines=true,
    showstringspaces=false,
    tabsize=4
}

\title{\textbf{Informe Técnico: Memory-Mapped I/O en MIPS32} \\
\large Análisis de Entrada/Salida Mapeada en Memoria}
\author{Arquitectura de Computadores}
\date{\today}

\begin{document}

\maketitle
\tableofcontents
\newpage

\section{Organización de la Memoria en Memory-Mapped I/O}

\subsection{Estructura de la Memoria}
En un sistema con memory-mapped I/O, el espacio de direcciones de memoria se divide en dos regiones principales:

\begin{itemize}
    \item \textbf{Región para RAM/RAM:} Espacio de direcciones asignado a la memoria principal (DRAM/SRAM)
    \item \textbf{Región para dispositivos I/O:} Espacio de direcciones reservado para los registros de los periféricos
\end{itemize}

\subsection{Región Típica de Mapeo}
En arquitecturas MIPS32 típicas, los dispositivos se mapean comúnmente en:
\begin{itemize}
    \item \texttt{0xFFFF0000 - 0xFFFFFFFF} (últimos 64KB del espacio de direcciones)
    \item Región \texttt{kseg1} en MIPS (direcciones no cacheables)
    \item Direcciones altas para evitar conflictos con la RAM convencional
\end{itemize}

\subsection{Implicaciones para lw y sw}
Las instrucciones \texttt{lw} (load word) y \texttt{sw} (store word) se utilizan tanto para acceder a memoria RAM como a registros de dispositivos:

\begin{lstlisting}[style=mipsstyle]
# Acceso a RAM
lw $t0, 0x1000($zero)     # Carga desde dirección de RAM

# Acceso a dispositivo I/O  
lw $t1, 0xFFFF0004($zero) # Carga desde registro de estado del sensor
sw $t2, 0xFFFF0000($zero) # Almacena en registro de control del sensor
\end{lstlisting}

El procesador no distingue entre ambos tipos de acceso; es el hardware externo quien interpreta la dirección y direcciona el acceso al dispositivo correspondiente.

\section{Memory-Mapped I/O vs. I/O por Puertos}

\subsection{Diferencias Principales}
\begin{table}[h]
\centering
\begin{tabular}{|p{5cm}|p{5cm}|}
\hline
\textbf{Memory-Mapped I/O} & \textbf{I/O por Puertos} \\
\hline
Mismo espacio de direcciones para memoria y E/S & Espacio de direcciones separado para E/S \\
\hline
Mismas instrucciones (lw/sw) para E/S y memoria & Instrucciones especiales (in/out) para E/S \\
\hline
Cualquier instrucción que acceda a memoria puede acceder a E/S & Sólo instrucciones específicas pueden acceder a E/S \\
\hline
\end{tabular}
\caption{Comparación de enfoques de E/S}
\end{table}

\subsection{Ventajas y Desventajas}

\textbf{Memory-Mapped I/O:}
\begin{itemize}
    \item \textbf{Ventajas:}
    \begin{itemize}
        \item No requiere instrucciones especiales
        \item Simplifica el diseño del procesador
        \item Mayor flexibilidad en el acceso a dispositivos
        \item Permite usar todas las instrucciones de memoria con dispositivos
    \end{itemize}
    \item \textbf{Desventajas:}
    \begin{itemize}
        \item Consume espacio de direcciones
        \item Puede requerir lógica adicional para decodificar direcciones
        \item Potenciales conflictos de caché
    \end{itemize}
\end{itemize}

\textbf{I/O por Puertos:}
\begin{itemize}
    \item \textbf{Ventajas:}
    \begin{itemize}
        \item Espacio de direcciones separado (no reduce memoria disponible)
        \item Instrucciones específicas permiten mejor control
        \item Más fácil implementar protección
    \end{itemize}
    \item \textbf{Desventajas:}
    \begin{itemize}
        \item Requiere instrucciones adicionales
        \item Complejidad añadida en el conjunto de instrucciones
        \item Menos flexibilidad en modos de acceso
    \end{itemize}
\end{itemize}

\subsection{¿Por qué MIPS32 usa Memory-Mapped I/O?}
MIPS32 utiliza principalmente memory-mapped I/O por varias razones:
\begin{enumerate}
    \item \textbf{Simplicidad RISC}: Mantiene el conjunto de instrucciones reducido
    \item \textbf{Eficiencia}: Aprovecha la jerarquía de memoria existente
    \item \textbf{Flexibilidad}: Permite usar modos de direccionamiento complejos
    \item \textbf{Estandarización}: Enfoque común en sistemas embebidos
\end{enumerate}

\section{Conflictos de Direcciones Solapadas}

\subsection{Problemas Potenciales}
Si dos dispositivos usan direcciones solapadas, pueden ocurrir:
\begin{itemize}
    \item Lecturas/escrituras dirigidas al dispositivo incorrecto
    \item Contención en el bus de datos
    \item Comportamiento indeterminado del sistema
    \item Posibles daños en el hardware
\end{itemize}

\subsection{Solución: Decodificación de Direcciones}
Para evitar conflictos, se implementa un decodificador de direcciones que:
\begin{enumerate}
    \item Asigna rangos exclusivos a cada dispositivo
    \item Utiliza señales de chip select (CS)
    \item Implementa jerarquía de decodificación
\end{enumerate}

\begin{lstlisting}[style=mipsstyle]
# Ejemplo de asignación exclusiva
.eqv SENSOR1_BASE 0xFFFF0000  # Rango: 0xFFFF0000-0xFFFF000F
.eqv SENSOR2_BASE 0xFFFF0010  # Rango: 0xFFFF0010-0xFFFF001F
.eqv SENSOR3_BASE 0xFFFF0020  # Rango: 0xFFFF0020-0xFFFF002F
\end{lstlisting}

\section{Simplificación del Conjunto de Instrucciones}

\subsection{Ventajas del Enfoque Unificado}
El memory-mapped I/O simplifica el diseño porque:
\begin{itemize}
    \item Reutiliza el mismo hardware de memoria
    \item No requiere nuevas instrucciones
    \item Mantiene la ortogonalidad del ISA
    \item Reduce la complejidad del decodificador de instrucciones
\end{itemize}

\subsection{Instrucciones Necesarias en I/O por Puertos}
Si se usara E/S por puertos, serían necesarias instrucciones adicionales como:
\begin{itemize}
    \item \texttt{in r1, port} - Leer del puerto
    \item \texttt{out port, r1} - Escribir al puerto
    \item \texttt{ins} - Entrada de cadena
    \item \texttt{outs} - Salida de cadena
\end{itemize}

\section{Operación a Nivel de Bus}

\subsection{Acceso a Dispositivos}
Cuando el procesador accede a una dirección de dispositivo:

\begin{enumerate}
    \item El procesador coloca la dirección en el bus de direcciones
    \item El decodificador de direcciones analiza la dirección
    \item Si la dirección corresponde a un dispositivo, activa la señal CS correspondiente
    \item El dispositivo responde en el bus de datos (lectura) o captura los datos (escritura)
    \item La memoria RAM permanece inactiva durante este acceso
\end{enumerate}

\subsection{Mecanismo de Selección}
\begin{verbatim}
Dirección: 0xFFFF0004
          ↓
Decodificador
    ↓           ↓
  RAM?        I/O?
  (No)        (Sí)
              ↓
        Activar CS_LUZ_ESTADO
              ↓
        Sensor responde en bus
\end{verbatim}

\section{Protección y Acceso de Programas sin Privilegios}

\subsection{Acceso no Autorizado}
Sí es posible que programas sin privilegios accedan a dispositivos si:
\begin{itemize}
    \item Las direcciones están en espacio de usuario
    \item No hay mecanismos de protección
    \item El sistema operativo no restringe el acceso
\end{itemize}

\subsection{Mecanismos de Protección}
Para evitar accesos no autorizados:

\begin{enumerate}
    \item \textbf{MMU (Memory Management Unit)}: Define permisos por página
    \item \textbf{Niveles de privilegio}: Solo modo kernel puede acceder a ciertas regiones
    \item \textbf{Sistema operativo}: Provee llamadas al sistema para acceso controlado
    \item \textbf{Trap handling}: Instrucciones ilegales generan excepciones
\end{enumerate}

En MIPS32 típicamente:
\begin{itemize}
    \item Dispositivos en \texttt{kseg1} (solo accesible en modo kernel)
    \item Excepción si usuario intenta acceder
    \item Drivers que operan en modo privilegiado
\end{itemize}

\section{Técnicas para Evitar Esperas Activas}

\subsection{Problema de la Espera Activa}
El código presentado en la práctica usa espera activa:
\begin{lstlisting}[style=mipsstyle]
EsperarListoLuz:
    lw $t0, LUZ_ESTADO
    beqz $t0, EsperarListoLuz  # Espera activa
\end{lstlisting}

\subsection{Técnicas Alternativas}

\textbf{1. Interrupciones:}
\begin{itemize}
    \item Dispositivo genera interrupción cuando está listo
    \item Procesador puede ejecutar otras tareas mientras espera
    \item Más eficiente para sistemas multitarea
\end{itemize}

\textbf{2. DMA (Direct Memory Access):}
\begin{itemize}
    \item Transferencia directa dispositivo-memoria
    \item Libera al procesador durante la transferencia
    \item Ideal para grandes volúmenes de datos
\end{itemize}

\textbf{3. Timeouts y Retries:}
\begin{itemize}
    \item Evitar bloqueos indefinidos
    \item Manejo robusto de errores
    \item Implementado en el sensor de presión
\end{itemize}

\begin{lstlisting}[style=mipsstyle]
EsperarConTimeout:
    li $t2, 1000000     # Contador de timeout
EsperarLoop:
    lw $t0, LUZ_ESTADO
    bnez $t0, Listo
    subi $t2, $t2, 1
    bnez $t2, EsperarLoop
    j TimeoutError      # Timeout alcanzado
\end{lstlisting}

\textbf{4. Planificación Cooperativa:}
\begin{itemize}
    \item Ceder el procesador mientras espera
    \item Usar llamadas al sistema como \texttt{yield()}
    \item Permite multiprogramación
\end{itemize}

\section{Análisis y Discusión de Resultados}

\subsection{Implementación de los Sensores}

\subsubsection{Sensor de Luminosidad}
La implementación sigue fielmente el modelo propuesto:
\begin{itemize}
    \item Inicialización explícita escribiendo 0x1 en \texttt{LuzControl}
    \item Verificación de estado antes de lectura
    \item Manejo de tres estados posibles (0, 1, -1)
    \item Retorno de código de estado junto con el valor
\end{itemize}

\subsubsection{Sensor de Presión}
Se implementó robustez adicional:
\begin{itemize}
    \item Detección de errores transitorios (-1)
    \item Reinicialización automática
    \item Reintento único como especificado
    \item Preservación de estado entre intentos
\end{itemize}

\subsubsection{Sensor de Tensión Arterial}
Caso más simple pero funcional:
\begin{itemize}
    \item Inicio de medición con comando 1
    \item Espera activa hasta completar
    \item Retorno de dos valores simultáneamente
\end{itemize}

\subsection{Eficiencia del Código}
\textbf{Ventajas:}
\begin{itemize}
    \item Código modular y reutilizable
    \item Manejo adecuado de errores
    \item Preservación de registros según convención
    \item Implementación clara y documentada
\end{itemize}

\textbf{Limitaciones:}
\begin{itemize}
    \item Uso de espera activa (busy-waiting)
    \item No hay manejo de timeouts
    \item Sin soporte para interrupciones
    \item Simulación simplificada
\end{itemize}

\subsection{Mejoras Propuestas}
Para una implementación más robusta:

\begin{enumerate}
    \item \textbf{Interrupciones:}
    \begin{lstlisting}[style=mipsstyle]
ConfigurarIntLuz:
    li $t0, 0x3        # Habilitar interrupción
    sw $t0, LUZ_CONTROL
    # ISR manejaría la lectura
    \end{lstlisting}
    
    \item \textbf{Buffering:}
    \begin{itemize}
        \item Almacenar múltiples lecturas
        \item Reducir frecuencia de acceso al hardware
    \end{itemize}
    
    \item \textbf{Timeouts configurables:}
    \begin{itemize}
        \item Evitar bloqueos permanentes
        \item Detectar hardware no responsive
    \end{itemize}
\end{enumerate}

\subsection{Conclusión}
El sistema implementado demuestra los principios fundamentales del memory-mapped I/O en MIPS32:
\begin{itemize}
    \item Acceso uniforme a memoria y dispositivos
    \item Mecanismos de sincronización mediante sondeo (polling)
    \item Manejo de errores y reintentos
    \item Importancia de la decodificación de direcciones
\end{itemize}

La principal limitación es la espera activa, que podría mejorarse con interrupciones en sistemas más sofisticados. Sin embargo, para sistemas embebidos simples o con tiempo real estricto, el enfoque de sondeo puede ser apropiado.

\end{document}